\section*{Chapitre \ref{ch:rc:sabr-in-rates-and-the-volatility-cube} --- SABR en taux et cube de volatilité}

\begin{solution}{exo:rc:sabr-in-rates-and-the-volatility-cube:1}
Non~: sa formule exige $F>0$ et $K>0$ (elle prend $\ln(F/K)$ et des puissances de $FK$).
Tout décalage supérieur à 0.50~\% rend $K+\zeta>0$ et $F+\zeta>0$~; en pratique on laisse une
marge, car la densité près du plancher se comporte mal.
\end{solution}

\begin{solution}{exo:rc:sabr-in-rates-and-the-volatility-cube:2}
La volatilité normale à la monnaie est inchangée ($\alpha$ multiplié par le même facteur)~:
le backbone est plat. Le strike à la monnaie d'hier se trouve désormais à 20 points de base
sous le forward, sa volatilité est donc celle du smile en $F-K=+20$ points
de base~: tout le smile s'est déplacé avec le forward.
\end{solution}

\begin{solution}{exo:rc:sabr-in-rates-and-the-volatility-cube:3}
Au sein d'une même section, $\beta$ et $\rho$ inclinent tous deux le smile et sont presque
interchangeables, de sorte qu'un calage des deux est instable d'un jour à l'autre~; on choisit
$\beta$ en fonction du backbone que le desk retient (la manière dont la volatilité à la monnaie
évolue avec le forward) et on le fixe pour la devise.
\end{solution}

\begin{solution}{exo:rc:sabr-in-rates-and-the-volatility-cube:4}
$25\times\bigl(1+\tfrac{2-3\times0.01}{24}\times0.36\times1\bigr) = 25\times1.02955 =
25.74$ points de base, la volatilité à la monnaie cotée.
\end{solution}

\begin{solution}{exo:rc:sabr-in-rates-and-the-volatility-cube:5}
Les cotations ont été générées par un modèle normal~; un \omterm{def:rc:sabr-in-rates-and-the-volatility-cube:shifted}{modèle lognormal décalé} avec un
faible décalage présente une volatilité fortement dépendante du niveau, puisque le forward décalé
$0.70~\%$ est petit, et sa forme de smile (courbure par rapport au skew) diffère
de celle du modèle normal, de sorte que trois paramètres ne peuvent pas reproduire exactement
six cotations. Un décalage plus grand fait que le \omterm{def:rc:sabr-in-rates-and-the-volatility-cube:shifted}{SABR décalé} se comporte davantage comme un modèle normal, et
le calage et l'aile concordent mieux.
\end{solution}

\begin{solution}{exo:rc:sabr-in-rates-and-the-volatility-cube:6}
Acheter le butterfly du modèle là où la densité est négative (position longue sur deux
swaptions receveuses au strike central, position courte d'une à chaque strike légèrement
au-dessus et au-dessous, autour de 0~\%)~: le modèle le price négativement, de sorte qu'un desk
valorisant avec le calage à 1~\% paierait pour le vendre. En pratique, ces strikes ne sont pas cotés, la fourchette cours acheteur--cours vendeur est large,
et les prix de marché de tels butterflies sont positifs~: l'arbitrage
se trouve dans le modèle, et un desk qui l'utiliserait serait victime de sélection adverse.
\end{solution}

\begin{solution}{exo:rc:sabr-in-rates-and-the-volatility-cube:7}
0.90~\%~: à 0.85~\% la densité du smile calibré est négative juste au-dessus du
plancher ($-0.80~\%$ à $-0.79~\%$), à 0.90~\% elle est positive de $-0.80~\%$ à $+1~\%$.
(Tout décalage doit dépasser 0.80~\% pour que les strikes cotés existent tout simplement.)
\end{solution}

\begin{solution}{exo:rc:sabr-in-rates-and-the-volatility-cube:8}
Les sections ont des forwards différents, de sorte que le même strike absolu se situe à
une moneyness différente dans chacune~; interpoler les volatilités strike par strike mélange
des points à la monnaie et des points d'aile, et le résultat n'a pas à être un smile SABR ni
exempt d'arbitrage. Interpoler les paramètres (ou les volatilités indexées sur la moneyness),
prendre le forward dans la courbe et vérifier la densité.
\end{solution}

\begin{solution}{pb:rc:sabr-in-rates-and-the-volatility-cube:1}
\textbf{1.} 60 points de base sous le forward, 10 sous la cotation la plus basse
($-0.80~\%$).
\textbf{2.} Sa formule exige des forwards et des strikes positifs.
\textbf{3.} \omterm{def:rc:sabr-in-rates-and-the-volatility-cube:shifted}{SABR décalé}, qui suppose que le taux reste au-dessus de $-\zeta$~; \omterm{def:rc:sabr-in-rates-and-the-volatility-cube:normal}{SABR normal},
qui laisse le taux prendre n'importe quelle valeur, avec des queues normales.
\textbf{4.} 0.77 point de base (décalage de 1~\%), 0.11 (3~\%), zéro pour le \omterm{def:rc:sabr-in-rates-and-the-volatility-cube:normal}{SABR normal} (le
modèle propre aux cotations).
\textbf{5.} \omterm{def:rc:sabr-in-rates-and-the-volatility-cube:normal}{SABR normal}, puis le décalage de 3~\%.
\textbf{6.} $0.177$, $0.280$ et $0.295$ points de base d'annuité, multipliés par $2.0\times
500$ millions~: 17\,717, 27\,999 et 29\,477~euros.
\textbf{7.} 10\,282~euros, soit 37~\% du prix à 3~\%.
\textbf{8.} Le taux modélisé ne peut pas descendre sous $-1~\%$, et sa volatilité diminue à mesure que
le taux s'approche du plancher, de sorte que la probabilité de finir sous $-0.90~\%$ est faible.
\textbf{9.} Le desk préférerait vendre au prix le plus élevé défendable, le client
acheter au plus bas~; le décalage de 1~\% est le modèle le plus favorable au client.
\textbf{10.} 0.90~\% (exercice~7).
\textbf{11.} La différence entre la dynamique des deux strikes~: la pente et
la courbure du smile sous les cotations, c'est-à-dire l'aile du modèle, plus l'écart entre les
strikes~; le desk est exposé à la réévaluation de l'aile par le marché.
\textbf{12.} Des taux proches de $-0.75~\%$ rendraient le plancher du décalage de 1~\% presque contraignant et
ses prix dénués de sens~; le desk aurait besoin d'un décalage plus grand, et les valorisations
sauteraient au moment où il le modifierait.
\textbf{13.} Son choix déplace les prix des strikes non cotés de plusieurs dizaines de pour cent
(ici 37~\%) tout en calant aussi bien les strikes cotés~; c'est de l'incertitude de
modèle, et une réserve couvre l'écart des prix entre les décalages admissibles.
\textbf{14.} Rien~: le taux est normal conditionnellement à la volatilité, non borné vers le bas,
ce qui accorde plus de poids aux taux très négatifs que n'importe quel modèle décalé.
\textbf{15.} Caler tous les modèles candidats~; vérifier les densités sur les strikes du portefeuille~;
comparer les prix d'aile entre modèles et avec tout point effectivement tradé (consultations de dealers)~;
faire un backtest du comportement du smile lorsque le forward a bougé.
\textbf{16.} Pour garder le plancher nettement sous tout taux plausible, de sorte que l'aile de
la distribution et les prix des strikes bas ne dépendent pas de l'endroit où se trouve le
plancher~; un petit décalage rend le smile otage de celui-ci.
\textbf{17.} Raisonnable au sein d'une devise s'il se situe nettement sous le \omterm{def:rc:sabr-in-rates-and-the-volatility-cube:elb}{plancher effectif
des taux}~; le choix fait partie de la convention de marché, qui rend les cotations
comparables.
\textbf{18.} Que les strikes cotés ne fixent le prix qu'au voisinage du forward, que trois
modèles standard donnent de 17\,700 à 29\,500~euros, et ce que le desk facture et pourquoi.
\textbf{19.} Résultat à retenir~: \emph{le strike moins-quatre-vingt-dix} coûte 17\,717~euros avec un décalage de 1~\%
et 27\,999~euros avec 3~\% (écart de 37~\%), et 0.90~\% est le plus petit décalage
qui garde la densité positive à partir de $-0.80~\%$.
\textbf{20.} L'hypothèse du modèle sur le niveau le plus bas que peuvent atteindre les taux et sur
leur volatilité à proximité de cette borne.
\end{solution}

\begin{solution}{iq:rc:sabr-in-rates-and-the-volatility-cube:1}
En renonçant aux cotations lognormales (Black), qui ne peuvent pas représenter de forwards
ou de strikes négatifs~: on cote des volatilités normales (Bachelier), ou des volatilités de Black d'un
taux décalé $F+\zeta$~; les modèles sont passés au \omterm{def:rc:sabr-in-rates-and-the-volatility-cube:normal}{SABR normal} ou au \omterm{def:rc:sabr-in-rates-and-the-volatility-cube:shifted}{SABR décalé}, puis plus tard
au SABR à frontière libre.
\iqlookfor{cotations normales et modèles décalés, et pourquoi Black échoue.}
\end{solution}

\begin{solution}{iq:rc:sabr-in-rates-and-the-volatility-cube:2}
$\beta$ fixe le backbone~: la manière dont la volatilité à la monnaie évolue avec la
forward ($\beta=0$ plat en termes normaux, $\beta=1$ plat en termes lognormaux) et une
partie du skew. Comme $\rho$ règle aussi le skew, caler les deux sur un même smile est
instable~; $\beta$ est fixé à partir de la vision du desk sur le backbone ou d'une
régression historique de la volatilité à la monnaie sur la forward.
\iqlookfor{backbone versus skew, et le problème d'identification.}
\end{solution}

\begin{solution}{iq:rc:sabr-in-rates-and-the-volatility-cube:3}
Pour chaque couple (échéance, tenor)~: forward et annuité issues de la multi-courbe~; nettoyer les
cotations (volatilités normales ou primes, straddles et strangles)~; fixer $\beta$ et le
décalage~; caler $\alpha,\rho,\nu$ avec bornes~; contrôler les erreurs et la régularité des
paramètres sur la grille~; interpoler les paramètres (variance en fonction de l'échéance)~;
vérifier la densité et l'arbitrage calendaire~; publier avec les diagnostics.
\iqlookfor{un pipeline avec contrôles, et pas seulement un calage.}
\end{solution}

\begin{solution}{iq:rc:sabr-in-rates-and-the-volatility-cube:4}
Normal SABR~: le smile se décale de dix points de base vers la droite, la volatilité normale à la
monnaie inchangée. Lognormal avec $\beta=1$~: la volatilité Black à la monnaie est inchangée, donc la
volatilité normale croît à peu près proportionnellement à la forward, et le smile se déplace en
termes relatifs (log-moneyness).
\iqlookfor{les deux backbones.}
\end{solution}

\begin{solution}{iq:rc:sabr-in-rates-and-the-volatility-cube:5}
Dans l'aile des bas strikes des échéances longues (grand $\nu^2T$), près de la borne inférieure
du taux, parce que le développement est asymptotique et représente mal la frontière absorbante.
Remèdes~: un décalage plus grand, un SABR exact ou numérique (EDP, intégration), le SABR à
frontière libre, ou une extrapolation des ailes avec une densité positive~; tester
systématiquement la densité sur les strikes du livre.
\iqlookfor{la cause et les remèdes.}
\end{solution}

\begin{solution}{iq:rc:sabr-in-rates-and-the-volatility-cube:6}
Regarder d'abord les données~: cotations périmées ou croisées, forward issue de la mauvaise
courbe, grille de strikes décalée par la forward~; puis le calage~: point de départ, bornes
atteintes ($\rho$ proche de $\pm1$, $\nu$ qui explose), échecs de conversion des cotations
profondes en volatilité du modèle~; comparer avec les paramètres de la veille et les sections
voisines~; ajouter une régularisation vers les voisins si la section est peu cotée.
\iqlookfor{les données avant l'optimiseur, et un contrôle systématique.}
\iqlookfor{les données avant l'optimiseur, et un contrôle systématique.}
\end{solution}
